\documentclass[11pt,letterpaper]{article}
\usepackage[margin=1in]{geometry}
\usepackage{hyperref}
\usepackage{booktabs}
\usepackage{longtable}
\usepackage{xcolor}
\usepackage{listings}
\usepackage{amsmath}
\usepackage{parskip}

\title{BVBRC-113 --- Independent Replication Report\\
\large Nakazono \textit{et al.} 2022 --- \textit{S. epidermidis} epidermin \& nukacin plasmids}
\author{Ollie (sub-agent), argo/argo:claude-opus-4.7 driver\\
LLM judge: argo:claude-sonnet-4.6}
\date{2026-07-05}

\begin{document}
\maketitle

\section{Paper Under Test}

\begin{itemize}
  \item \textbf{Citation:} Nakazono K, Le MN, Kawada-Matsuo M, Kimheang N, Hisatsune J, Oogai Y, Nakata M, Nakamura N, Sugai M, Komatsuzawa H. \textit{Complete sequences of epidermin and nukacin encoding plasmids from oral-derived Staphylococcus epidermidis and their antibacterial activity.} PLoS ONE 17(1):e0258283, 18 Jan 2022.
  \item \textbf{PMID:} 35041663 \quad \textbf{PMC:} PMC8765612 \quad \textbf{DOI:} 10.1371/journal.pone.0258283
  \item \textbf{License:} CC-BY 4.0
  \item \textbf{Deposited data:} NCBI \texttt{OK031036} (pEpi56, 64{,}386 bp) and \texttt{OK031035} (pNuk650, 26{,}160 bp).
  \item \textbf{Wave:} BVBRC top-up 2026-07-05.
\end{itemize}

\section{Summary of Paper}

Nakazono \textit{et al.} screened 150 \textit{S. epidermidis} oral isolates from 287 volunteers for bacteriocin activity against a hypersusceptible \textit{S. aureus} MW2 \textit{braRS} indicator. Two producers, \textbf{KSE56} and \textbf{KSE650}, were characterised. Illumina + MinION hybrid assembly (Unicycler v0.4.8, RAST annotation) produced complete genomes; the bacteriocin-encoding plasmids \textbf{pEpi56} (epidermin) and \textbf{pNuk650} (nukacin) were deposited under OK031036/OK031035. The paper additionally reports MS-verified purification, susceptibility spectra, and \textit{M.\ luteus} co-culture experiments.

Eight testable in-silico claims (C1--C8) are extracted from Tables~2--3 and Figures~2--3 and re-checked against the deposited public sequences.

\section{Method}

All computation executed 2026-07-05 on CherryRd in a local Python venv (Python 3.14.6, Biopython 1.87). LLM judge via free localhost Argo proxy (\texttt{127.0.0.1:44497}).

\begin{enumerate}
  \item \textbf{Paper acquisition:} \texttt{esummary/efetch} $\rightarrow$ \texttt{work/paper.xml} (225{,}843 B).
  \item \textbf{Sequence acquisition:} GenBank + FASTA fetched for \texttt{OK031036 OK031035 KP702950 X62386 U77778} via NCBI E-utilities.
  \item \textbf{Structural verification} (\texttt{analyze\_plasmids.py}): Biopython \texttt{SeqIO} parsing $\rightarrow$ length, topology, CDS counts, \textit{epi*}/\textit{nuk*} gene inventory.
  \item \textbf{Bacteriocin identity} (\texttt{bacteriocin\_align.py}): direct position-by-position mismatch counting for \textit{epiA} (KSE56 vs T\"u3298 X62386) and \textit{nukA} (KSE650 vs IVK45 KP702950); mature-peptide comparison over terminal 27 aa.
  \item \textbf{Comparative ORF-delta} (\texttt{compare\_plasmids.py}): \texttt{blastp} of all 29 pNuk650 CDS against a pIVK45 (17 CDS) proteome db; ``true ortholog'' called at $\text{pident}\geq 30\%$ AND $\text{qcov}\geq 50\%$.
  \item \textbf{LLM-judge verdict} (\texttt{llm\_judge.py}): three JSON evidence blobs + numbered claims table $\rightarrow$ argo:claude-sonnet-4.6 (opus-4.7/4.8 both 502'd on $\sim$30\,kB payload).
\end{enumerate}

\section{Results}

\subsection{Structural plasmid claims}

\begin{longtable}{lllc}
\toprule
Claim & Paper & Independent measurement & Verdict \\
\midrule
pEpi56 length & 64{,}386 bp & 64{,}386 bp (OK031036) & MATCH \\
pEpi56 topology & circular & circular & MATCH \\
pEpi56 ORFs & 81 & 81 CDS & MATCH \\
pNuk650 length & 26{,}160 bp & 26{,}160 bp (OK031035) & MATCH \\
pNuk650 topology & circular & circular & MATCH \\
pNuk650 ORFs & 29 & 29 CDS & MATCH \\
pIVK45 length & 21{,}840 bp & 21{,}840 bp (KP702950) & MATCH \\
pNuk650 $>$ pIVK45 & yes ($\sim$+4.3 kb) & +4{,}320 bp & MATCH \\
pNuk650 extra ORFs & +7 & +12 raw / +13 no-ortholog & PARTIAL \\
\bottomrule
\end{longtable}

\subsection{Epidermin (KSE56 vs T\"u3298 X62386)}

Both \textit{epiA} CDS are 159 nt / 52 aa.

\begin{itemize}
  \item Nucleotide mismatches: \textbf{2} (positions 34 and 152, both synonymous)
  \item Amino-acid mismatches: \textbf{0} (100.0\% aa identity)
\end{itemize}

Paper claim (2 nt / 100\% aa) $\Rightarrow$ \textbf{EXACT MATCH}.

\subsection{Nukacin (KSE650 vs IVK45 KP702950)}

Both \textit{nukA} CDS are 174 nt / 57 aa.

\begin{itemize}
  \item Prepeptide aa mismatches: \textbf{1} (position 4, L$\leftrightarrow$F)
  \item Mature peptide (terminal 27 aa): identical
\end{itemize}

Paper claim (1 prepeptide mm at position 4; mature peptides identical) $\Rightarrow$ \textbf{EXACT MATCH}.

\subsection{Bacteriocin gene clusters}

\begin{itemize}
  \item pEpi56: \texttt{epiP epiQ epiD epiC epiB epiA epiT' epiH epiF epiE epiG} --- 11 named \textit{epi} loci; no explicit \texttt{epiY} annotation.
  \item pNuk650: \texttt{nukA nukM nukT nukF nukE nukG nukH} --- all 7 nukacin genes present.
\end{itemize}

\subsection{LLM-judge structured verdict}

Overall verdict: \textbf{PARTIAL}. Overall score: \textbf{74 / 100}. Six of eight claims MATCH; two are PARTIAL (the ``+7 ORFs'' count and the \textit{epi} cluster naming). No claim contradicted at the sequence level.

\section{GENUINE CRITIQUE}

This section is a candid appraisal of the replication as a whole and of the paper on the merits of what did and did not reproduce.

\subsection{What genuinely replicated (and why that is meaningful)}

The two most quantitatively specific molecular claims in the paper reproduce \emph{to the nucleotide}:
\begin{itemize}
  \item \textit{epiA} (KSE56 vs T\"u3298): 2 nt mismatches, 0 aa mismatches, 100.0\% aa identity.
  \item \textit{nukA} (KSE650 vs IVK45): 1 aa mismatch at prepeptide position 4 (Leu$\leftrightarrow$Phe); mature peptide entirely identical over its terminal 27 residues.
\end{itemize}
These are the kinds of narrow, falsifiable claims that would immediately shatter under sloppy sequencing or misannotation. The fact that both survived a from-scratch re-derivation against the primary deposits is strong evidence that the sequencing, base-calling, and CDS annotation reported in the paper are trustworthy at the locus of interest --- and by extension, that the paper's identification of KSE56 and KSE650 as bona-fide epidermin/nukacin producers rests on a solid genomic foundation. Plasmid lengths (64{,}386 bp; 26{,}160 bp), topologies (both circular), and total CDS counts (81; 29) all recovered exactly from the deposited GenBank records, confirming the finished, complete-assembly nature of the deposits.

\subsection{What did NOT replicate cleanly}

Two claims came up PARTIAL:
\begin{enumerate}
  \item \textbf{``pNuk650 has an additional seven ORFs vs pIVK45.''} The raw CDS-count delta between OK031035 (29 CDS) and KP702950 (17 CDS) is \textbf{+12}, not +7; a BLASTP-based no-ortholog count is \textbf{+13}. Under no straightforward counting method against the deposited pIVK45 annotation does the number ``7'' emerge. The most benign explanations --- (a) annotation-depth asymmetry (pIVK45 is more coarsely annotated in KP702950 than pNuk650), and (b) a narrower definition scoped only to the $\sim$8 kbp inserted region flanking the shared nukacin cluster --- are plausible but neither is verifiable without the authors' custom re-annotation. This is a soft factual slip in the abstract, not a structural error; the core qualitative claim (pNuk650 is larger and carries extra content) is correct.
  \item \textbf{Epidermin cluster naming (\textit{epiABCDEFGHPQTY}).} The deposited record contains named annotations for 11 \textit{epi} loci but no explicit \texttt{epiY}. Since pEpi56 carries an intact \textit{epiT} (per the paper's own finding), \textit{epiY} presence-vs-absence is a comparative statement about T\"u3298, not a structural feature of pEpi56 --- so this is a naming/notation ambiguity rather than a substantive discrepancy. Still, an in-silico verifier that reads only the deposited record cannot enumerate \textit{epiY} from OK031036 alone.
\end{enumerate}

\subsection{What was categorically out of reach}

An honest replication has to state its own boundary. The following wet-lab claims are simply \textbf{NOT TESTED} by this work and cannot be tested from public sequence data alone:
\begin{itemize}
  \item Bacteriocin purification (cation-exchange + HPLC), ESI-MS mass verification of the mature lantibiotics.
  \item Plasmid curing with acriflavine and the loss-of-activity phenotype.
  \item MW2 \textit{braRS} susceptibility assays (the very hypersusceptible indicator that drove the original screen).
  \item \textit{M.\ luteus} co-culture qPCR and antibacterial spectra against $\geq 15$ oral/skin commensals.
\end{itemize}
These wet claims constitute the biological punchline of the paper. A PARTIAL verdict is honest precisely because it does not paper over the fact that the paper's activity-and-spectrum story is un-audited by this analysis.

\subsection{Methodological caveats on our own side}

\begin{itemize}
  \item \textbf{Ortholog thresholds are arbitrary.} We used pident $\geq 30\%$ and qcov $\geq 50\%$; sliding these could easily shift the ``no-ortholog'' count by $\pm$2. This is precisely why we do not claim the paper is \textit{wrong} about +7 --- only that we cannot recover +7 under any obvious rule.
  \item \textbf{Annotation, not re-annotation.} We took RAST/GenBank CDS calls as given on both sides; a normalized re-annotation (e.g., Prokka or Bakta on both plasmids) would be a fairer basis for a count-delta claim.
  \item \textbf{LLM judge is a formalism, not an oracle.} The 74/100 score is a structured aggregation of eight per-claim verdicts, not an independent semantic judgment; opus-4.7 and opus-4.8 both 502'd on the payload, and only sonnet-4.6 completed --- the score should be read as a rubric summary, not as a model consensus.
  \item \textbf{Marker OCR absent.} The extraction pipeline's Marker output was not present at write time; all quantitative claims here are anchored on the deposited sequences and the paper text, not on any parsed table extraction.
\end{itemize}

\subsection{Bottom-line critique}

The paper is on \emph{sound sequence-analysis footing}. Every sequence-level claim we could verify verified; the one numeric discrepancy (``+7 ORFs'') is peripheral to the biological argument and is almost certainly attributable to differing annotation conventions rather than to an error in the underlying assembly. The paper's core biology --- that KSE56 and KSE650 encode bona-fide epidermin and nukacin on complete circular plasmids, with a mature nukacin peptide identical to that of IVK45 --- is credible and reproducibly documented in the public record. The paper's activity-spectrum and MS-purification claims remain unverified by this replication and should not be misrepresented as replicated.

\section{Verdict}

\textbf{PARTIAL --- 74/100.} Six of eight in-silico claims MATCH exactly; two are PARTIAL (numeric ``+7'' ORFs; \textit{epi} cluster naming). No claim was contradicted at the sequence level. Wet-lab claims are out of scope.

\textbf{Recommendation:} Solid PARTIAL. The paper is trustworthy at the sequence-and-structure level; the ``+7 ORFs'' figure should be treated as approximate and re-derived from a normalized re-annotation before being used elsewhere.

\section{Reproducibility Bundle}

\begin{itemize}
  \item \textbf{Code:} \texttt{work/analyze\_plasmids.py}, \texttt{work/bacteriocin\_align.py}, \texttt{work/compare\_plasmids.py}, \texttt{work/llm\_judge.py}.
  \item \textbf{Data:} \texttt{work/sequences/\{OK031036,OK031035,KP702950,X62386,U77778\}.\{gb,fasta\}}, \texttt{work/paper.xml}.
  \item \textbf{Evidence:} \texttt{report/evidence/\{plasmid\_summary,bacteriocin\_alignment,pNuk650\_vs\_pIVK45\_blast\}.json}, \texttt{report/evidence/llm\_judge\_verdict.txt}.
  \item \textbf{Re-run:} \texttt{cd work \&\& python3 -m venv .venv \&\& source .venv/bin/activate \&\& pip install biopython \&\& python analyze\_plasmids.py \&\& python bacteriocin\_align.py \&\& python compare\_plasmids.py \&\& python llm\_judge.py}. Wall time $\approx$ 30 s.
\end{itemize}

\end{document}
